\documentclass[11pt]{article}

\usepackage[T1]{fontenc}
\usepackage[utf8]{inputenc}
\usepackage{lmodern}
\usepackage[margin=1in]{geometry}
\usepackage{amsmath,amssymb,amsthm,mathtools}
\usepackage{booktabs}
\usepackage{array}
\usepackage{enumitem}
\usepackage{microtype}
\usepackage{xcolor}
\usepackage{hyperref}
\usepackage{fancyhdr}

\definecolor{linkblue}{RGB}{31,78,121}
\definecolor{resultgreen}{RGB}{32,112,72}
\hypersetup{
  colorlinks=true,
  linkcolor=linkblue,
  urlcolor=linkblue,
  citecolor=linkblue,
  pdftitle={An 18-Vertex Counterexample to WOWII Conjecture 59},
  pdfauthor={Prepared for reproducible verification},
  pdfsubject={Graph-theoretic counterexample with exhaustive verification},
  pdfkeywords={graph theory, Havel-Hakimi residue, induced forest, bipartite subgraph, counterexample}
}

\pagestyle{fancy}
\fancyhf{}
\lhead{WOWII Conjecture 59}
\rhead{18-vertex counterexample}
\cfoot{\thepage}
\setlength{\headheight}{14pt}

\newtheorem{theorem}{Theorem}
\newtheorem{lemma}{Lemma}
\newtheorem{proposition}{Proposition}

\newcommand{\res}{\operatorname{residue}}
\newcommand{\ceil}[1]{\left\lceil #1 \right\rceil}
\newcommand{\floor}[1]{\left\lfloor #1 \right\rfloor}
\newcommand{\set}[1]{\left\{#1\right\}}

\title{\textbf{An 18-Vertex Counterexample\\to WOWII Conjecture 59}}
\author{A reproducible counterexample note}
\date{July 23, 2026}

\begin{document}
\maketitle

\begin{abstract}
We give an explicit connected simple graph $G$ on $18$ vertices for which
\[
  \res(G)=10,\qquad b(G)=17,\qquad f(G)=13,
\]
where $b(G)$ is the maximum order of an induced bipartite subgraph and $f(G)$
is the maximum order of an induced forest. Consequently,
\[
  \ceil{\sqrt{\res(G)b(G)}}=\ceil{\sqrt{170}}=14>13=f(G),
\]
disproving the stated form of Written on the Wall II, Conjecture 59. The
argument below proves all three invariant values exactly. An accompanying
standard-library Python program independently enumerates every one of the
$2^{18}$ induced vertex subsets.
\end{abstract}

\section{The conjectured inequality}

The form of WOWII Conjecture 59 considered here is
\begin{equation}\label{eq:conjecture}
  f(G)\ \ge\ \ceil{\sqrt{\res(G)b(G)}}
\end{equation}
for every finite simple connected graph $G$. Here:
\begin{itemize}[leftmargin=2em]
  \item $f(G)$ is the maximum number of vertices in an induced forest of $G$;
  \item $b(G)$ is the maximum number of vertices in an induced bipartite
        subgraph of $G$; and
  \item $\res(G)$ is the Havel--Hakimi residue, namely the number of zeros
        left when the degree sequence is reduced until all remaining entries
        are zero.
\end{itemize}

\begin{theorem}\label{thm:main}
There exists a connected graph $G$ on $18$ vertices with
\[
\res(G)=10,\qquad b(G)=17,\qquad f(G)=13.
\]
In particular, $G$ violates \eqref{eq:conjecture}.
\end{theorem}

\section{Construction of the graph}

Let
\[
 A=\set{a_0,a_1,a_2,a_3,a_4},\qquad
 B=\set{b_0,b_1,b_2,b_3,b_4},
\]
let $h$ be one additional vertex, and let
\[
 L=\set{\ell_1,\ell_2,\ldots,\ell_7}
\]
be seven further vertices. Thus
\[
 V(G)=A\cup B\cup\set{h}\cup L.
\]

The graph induced by $A\cup B$ is bipartite. Its neighborhoods on the $B$ side
are given in Table~\ref{tab:neighborhoods}.

\begin{table}[ht]
\centering
\caption{Neighborhoods in the ten-vertex bipartite core.}
\label{tab:neighborhoods}
\begin{tabular}{@{}c l c@{}}
\toprule
Vertex & Neighborhood in $A$ & Degree in the core \\
\midrule
$b_0$ & $\set{a_0,a_1}$                         & $2$ \\
$b_1$ & $\set{a_0,a_1,a_2}$                     & $3$ \\
$b_2$ & $\set{a_0,a_1,a_2,a_3}$                 & $4$ \\
$b_3$ & $A$                                      & $5$ \\
$b_4$ & $A$                                      & $5$ \\
\bottomrule
\end{tabular}
\end{table}

Finally, join $h$ to every other vertex and add no further edges. In
particular, each $\ell_i$ is a pendant vertex adjacent only to $h$. Since $h$
is universal, $G$ is connected. The core has $19$ edges and the universal hub
adds $17$, so $|E(G)|=36$.

For machine verification we use the numerical labeling
\[
 a_i=i,\qquad b_i=5+i,\qquad h=10,\qquad
 L=\set{11,12,13,14,15,16,17}.
\]

\section{The induced bipartite number}

\begin{proposition}\label{prop:b}
The graph $G$ satisfies $b(G)=17$.
\end{proposition}

\begin{proof}
Deleting $h$ leaves the bipartite core on $A\cup B$ together with seven
isolated vertices. Hence $G-h$ is an induced bipartite subgraph on $17$
vertices, and therefore $b(G)\ge17$.

On the other hand, $a_0b_0$, $b_0h$, and $ha_0$ are all edges. Thus
$\set{h,a_0,b_0}$ spans a triangle, so the full $18$-vertex graph is not
bipartite. No induced bipartite subgraph can therefore have order $18$.
Hence $b(G)=17$.
\end{proof}

\section{The induced forest number}

Write $H=G[A\cup B]$ for the ten-vertex bipartite core.

\begin{lemma}\label{lem:alpha-core}
The independence number of $H$ is $5$.
\end{lemma}

\begin{proof}
The part $A$ is independent, so $\alpha(H)\ge5$. The five edges
\[
 a_0b_0,\quad a_1b_1,\quad a_2b_2,\quad a_3b_3,\quad a_4b_4
\]
form a perfect matching of $H$. An independent set contains at most one
endpoint from each matching edge, so it has size at most $5$. Therefore
$\alpha(H)=5$.
\end{proof}

\begin{lemma}\label{lem:seven-core}
Every set of seven vertices of $H$ contains a $4$-cycle.
\end{lemma}

\begin{proof}
Let $S\subseteq V(H)$ have $|S|=7$, and put
\[
 a=|S\cap A|,\qquad b=|S\cap B|.
\]
Since $a+b=7$ and $a,b\le5$, the possibilities are
\[
 (a,b)\in\set{(5,2),(4,3),(3,4),(2,5)}.
\]
The neighborhoods in Table~\ref{tab:neighborhoods} are nested:
\[
 N(b_0)\subset N(b_1)\subset N(b_2)\subset N(b_3)=N(b_4)=A,
\]
with sizes $2,3,4,5,5$.

\begin{enumerate}[label=(\alph*),leftmargin=2.5em]
  \item If $(a,b)=(5,2)$, the two selected vertices of $B$ have at least two
        common neighbors in $A$.
  \item If $(a,b)=(4,3)$, choose the two selected $B$-vertices with largest
        neighborhoods. Their common neighborhood has size at least $3$.
        Since only one vertex of $A$ is omitted, at least two of those common
        neighbors lie in $S$.
  \item If $(a,b)=(3,4)$, the two selected $B$-vertices with largest
        neighborhoods have at least four common neighbors in $A$. The
        complement of that common neighborhood inside $A$ has size at most
        one, so any selected three vertices of $A$ include at least two common
        neighbors.
  \item If $(a,b)=(2,5)$, both $b_3$ and $b_4$ are selected and are adjacent
        to both selected vertices of $A$.
\end{enumerate}

In every case, two selected vertices of $B$ have two common selected neighbors
in $A$. These four vertices span a $K_{2,2}$ and hence a $4$-cycle.
\end{proof}

\begin{proposition}\label{prop:f}
The graph $G$ satisfies $f(G)=13$.
\end{proposition}

\begin{proof}
The set
\[
 B\cup\set{h}\cup L
\]
has $5+1+7=13$ vertices and induces a star centered at $h$. Therefore
$f(G)\ge13$.

For the reverse inequality, let $S\subseteq V(G)$ have $|S|=14$.

If $h\in S$, then at most seven vertices of $S\setminus\set{h}$ belong to
$L$, so $S$ contains at least six vertices of $H$. By
Lemma~\ref{lem:alpha-core}, those six core vertices contain an edge $xy$.
The universal hub is adjacent to both $x$ and $y$, so $hxyh$ is a triangle in
$G[S]$. Thus $G[S]$ is not a forest.

If $h\notin S$, then at most seven vertices of $S$ belong to $L$, so $S$
contains at least seven vertices of $H$. By Lemma~\ref{lem:seven-core}, those
core vertices contain a $4$-cycle. Again $G[S]$ is not a forest.

No induced forest has $14$ vertices, so $f(G)\le13$. Together with the lower
bound, this gives $f(G)=13$.
\end{proof}

\section{The Havel--Hakimi residue}

The degrees are
\begin{itemize}[leftmargin=2em]
  \item $17$ for the hub $h$;
  \item $6,6,5,4,3$ on $A$;
  \item $3,4,5,6,6$ on $B$; and
  \item $1$ for each of the seven leaves.
\end{itemize}
Thus the sorted degree sequence is
\[
 d_0=[17,6^4,5^2,4^2,3^2,1^7],
\]
where $x^k$ denotes $k$ copies of $x$. Successive Havel--Hakimi reductions are
\begin{align*}
 d_0&=[17,6^4,5^2,4^2,3^2,1^7],\\
 d_1&=[5^4,4^2,3^2,2^2,0^7],\\
 d_2&=[4^3,3^4,2^2,0^7],\\
 d_3&=[3^4,2^4,0^7],\\
 d_4&=[2^7,0^7],\\
 d_5&=[2^4,1^2,0^7],\\
 d_6&=[2,1^4,0^7],\\
 d_7&=[1^2,0^9],\\
 d_8&=[0^{10}].
\end{align*}
The terminal list has ten zeros, and hence
\begin{equation}\label{eq:residue}
  \res(G)=10.
\end{equation}

\section{Contradiction to Conjecture 59}

Propositions~\ref{prop:b} and~\ref{prop:f}, together with
\eqref{eq:residue}, give
\[
 \res(G)b(G)=10\cdot17=170=13^2+1.
\]
Therefore
\[
 \ceil{\sqrt{\res(G)b(G)}}=\ceil{\sqrt{170}}=14,
\]
whereas $f(G)=13$. The conjectured inequality reads
\[
 13\ge14,
\]
which is false. This proves Theorem~\ref{thm:main}.

\begin{center}
\setlength{\fboxsep}{10pt}
\color{resultgreen}
\fbox{\color{black}\parbox{0.82\linewidth}{\centering
\textbf{Conclusion.} The graph above is a connected 18-vertex counterexample
to the stated form of WOWII Conjecture 59.}}
\end{center}

\section{Independent exhaustive verification}

The repository includes
\begin{center}
\texttt{verification/verify\_counterexample.py}.
\end{center}
The program uses only the Python standard library. It constructs the exact
edge set above, checks connectivity, computes the Havel--Hakimi residue, and
examines all $2^{18}=262{,}144$ vertex subsets. Each induced subgraph is tested
for bipartiteness by breadth-first two-coloring and for acyclicity by a
union--find cycle test. It then asserts the exact values
\[
 \res(G)=10,\qquad b(G)=17,\qquad f(G)=13,
\]
and the final strict inequality $13<14$.

From the repository root, run
\begin{verbatim}
python3 verification/verify_counterexample.py
\end{verbatim}
The expected output is
\begin{verbatim}
|V| = 18, |E| = 36, connected = True
degree sequence = [17, 6, 6, 6, 6, 5, 5, 4, 4, 3, 3, 1, 1, 1, 1, 1, 1, 1]
Havel--Hakimi residue = 10
b(G) = 17; witness = [0, 1, 2, 3, 4, 5, 6, 7, 8, 9, 11, 12, 13, 14, 15, 16, 17]
f(G) = 13; witness = [0, 1, 2, 3, 4, 5, 11, 12, 13, 14, 15, 16, 17]
ceil(sqrt(residue*b)) = ceil(sqrt(170)) = 14
Conjectured inequality: 13 >= 14  ->  False
\end{verbatim}

\section{Sources, scope, and attribution}

The original WOWII list is available at
\begin{center}
\url{http://cms.dt.uh.edu/faculty/delavinae/research/wowII/all.html#conj59}.
\end{center}
The corresponding Formal Conjectures transcription is at
\begin{center}
\url{https://github.com/google-deepmind/formal-conjectures/blob/main/FormalConjectures/WrittenOnTheWallII/GraphConjecture59.lean}.
\end{center}
A related upstream Lean-certification pull request is
\begin{center}
\url{https://github.com/google-deepmind/formal-conjectures/pull/4583}.
\end{center}
This note concerns the exact inequality and invariant meanings stated in
Section~1. It makes no claim of historical or mathematical priority; related
upstream work should be consulted when assigning attribution.

\appendix
\section{Explicit numeric edge list}

Under the labeling
\[
 A=\set{0,1,2,3,4},\quad B=\set{5,6,7,8,9},\quad h=10,
 \quad L=\set{11,12,13,14,15,16,17},
\]
the core edges are
\begin{align*}
 &(0,5),(1,5),\\
 &(0,6),(1,6),(2,6),\\
 &(0,7),(1,7),(2,7),(3,7),\\
 &(0,8),(1,8),(2,8),(3,8),(4,8),\\
 &(0,9),(1,9),(2,9),(3,9),(4,9).
\end{align*}
In addition, $(10,v)$ is an edge for every $v\ne10$. There are no other
edges.

\end{document}
